\documentclass[../../main]{subfiles}

\begin{document}

本章では\c LaTeXにおける基本的な環境や機能の使い方について確認します\p 
レポートクラス(\code{jsreport}など)では\c 文書の最上位の区切りとして\code{\textbackslash chapter}を使用します\p これにより\c 節(\code{\textbackslash section})の番号が「章番号.節番号」のように階層化されて表示されます\p 

\section{各種コマンドと環境のテスト}
ここでは\c よく使われるコマンドや環境(リスト\c 図の挿入\c コード引用\c フレームなど)の記述例を示します\p 

\subsection{箇条書き(リスト)環境}
基本の箇条書き環境です\p 

\begin{itemize}
  \item 項目1
  \item 項目2
    \begin{enumerate}
      \item 番号付きサブ項目1
      \item 番号付きサブ項目2
    \end{enumerate}
  \item 項目3
\end{itemize}

\subsection{画像の挿入}
\code{graphicx}パッケージを用いて画像ファイルを挿入する例です(図\ref{fig:sample})\p 

\begin{figure}[htbp]
  \centering
  \includegraphics[width=5cm]{figures/img01.jpeg}
  \caption{サンプルの図}
  \label{fig:sample}
\end{figure}

\subsection{ソースコードの挿入}
\code{listings}パッケージを用いたソースコードの挿入例です(コード\ref{src:python_hello})\p 

\begin{lstlisting}[language=Python, caption=PythonのHello World, label=src:python_hello]
def hello_world():
    print("Hello, LaTeX!")

if __name__ == "__main__":
    hello_world()
\end{lstlisting}

\subsection{ハイパーリンク}
\code{hyperref}パッケージを用いたリンク挿入の例です\p 
直接URLを記述する場合は \url{https://www.google.com} のように記述します\p 
テキストにリンクを埋め込む場合は \href{https://www.google.com}{こちら} のように記述します\p 

\subsection{囲み枠(フレーム)環境}
\code{framed}パッケージを用いたフレーム環境の例です\p 

\begin{framed}
これはフレーム環境の中に書かれた文章です\p 
重要な文章や注意書きを強調するのに便利です\p 
\cite{einstein1905}
\end{framed}

\end{document}
